\documentclass[10pt,a4paper,twocolumn]{article}

% Chinese language support
\usepackage[UTF8,fontset=none]{ctex}

% Set Traditional Chinese for figure, table and reference captions
\renewcommand{\figurename}{圖}
\renewcommand{\tablename}{表}
\renewcommand{\refname}{參考文獻}

% Font customization (load after ctex)  
\usepackage{fontspec}

% Chinese font - Using Kaiti (楷體) with full path
% Kaiti is located in AssetsV2 directory (macOS dynamic fonts)
\setCJKmainfont[Path=/System/Library/AssetsV2/com_apple_MobileAsset_Font8/88d6cc32a907955efa1d014207889413890573be.asset/AssetData/]{Kaiti.ttc}
\setCJKsansfont{Heiti TC Light}
\setCJKmonofont{Heiti TC Light}

% Other packages
\usepackage{amsmath}
\usepackage{amssymb}
\usepackage{graphicx}
\usepackage{hyperref}
\usepackage{geometry}
\usepackage{setspace}
\usepackage{enumitem}
\usepackage{booktabs}
\usepackage{multirow}
\usepackage{makecell}
\usepackage{array}
\usepackage{xcolor}
\usepackage{colortbl}

% Page settings - tighter for conference format
\geometry{left=1.8cm,right=1.8cm,top=2cm,bottom=2cm,columnsep=0.6cm}

% Line spacing and paragraph settings - compressed
\linespread{1.1}  % Reduced from 1.25 for space efficiency
\setlength{\parskip}{0pt}
\setlength{\parindent}{2em}

% Avoid excessive vertical stretching
\raggedbottom

% Compact lists to save vertical space
\setlist{nosep}

% Title format settings
\usepackage{titling}
\pretitle{\begin{center}\fontsize{24}{28}\selectfont}
\posttitle{\par\end{center}\vskip 0.5em}
\preauthor{\begin{center}\fontsize{12}{14}\selectfont}
\postauthor{\end{center}}
\predate{\begin{center}\fontsize{12}{14}\selectfont}
\postdate{\par\end{center}}

% Title information
\title{基於 Kolmogorov-Arnold 圖神經網路之可解釋且穩健的微服務根因分析}

% Author information
\author{
    \textbf{黃冠澔} \quad \textbf{賴瑀辰} \quad \textbf{賴冠州} \\[0.5em]
    國立臺中教育大學資訊工程學系 \\[0.3em]
    \texttt{\{bcs113116, bcs113115\}@gm.ntcu.edu.tw} \quad \texttt{kclai@mail.ntcu.edu.tw}
}
\date{}

\begin{document}

% Title and abstract in two-column format
\maketitle

% Abstract section
\begin{abstract}
\noindent
儘管圖神經網路（GNN）在微服務根因分析中展現卓越的特徵提取能力，其黑箱決策機制與缺乏領域歸納偏差的特性，嚴重阻礙其在\textbf{具高風險（High-stakes，即系統故障可能導致重大損失）之關鍵維運場景}中的應用。為解決\textbf{表達力與可解釋性的兩難}，本研究提出 \textbf{GNN\_KAN}，一種基於 Kolmogorov–Arnold 表示定理的創新架構。GNN\_KAN 將多變量函數分解為可視化的一維基底函數組合，在保留深度學習強大擬合能力的同時，賦予模型\textbf{內在可解釋性（Intrinsic Interpretability）}，為根因分析提供了從隱式特徵映射轉向顯式函數表達的新視角。本研究進一步建構一套領域感知（Domain-aware）的特徵工程框架，系統性排除了特徵中的後驗洩漏（Posterior Leakage）風險。實驗結果顯示，GNN\_KAN 在多個包含多源遙測資料的核心資料集中，於 Top-k 準確度與 Precision@k 等關鍵指標上均優於現有統計基線方法（如 BARO）。在複雜拓樸場景（Train Ticket）中，GNN\_KAN 的 Precision@1 達到 0.256，相較基線方法提升 88\%。此外，本研究基於廣泛實證研究，揭示不同基底函數與系統拓樸複雜度間的內在映射關係，為實務應用提供系統化選擇策略。
\end{abstract}

\noindent\textbf{關鍵詞：}根因分析、圖神經網路、Kolmogorov-Arnold Networks、微服務系統、可解釋性人工智慧

\vspace{0.5em}

% 1. Introduction
\section{緒論}

\subsection{研究背景與動機}

隨著微服務架構普及 \cite{microservices_survey}，系統拓樸日益複雜，傳統根因分析（RCA）面臨嚴峻挑戰。當系統發生異常時，上游服務異常迅速沿依賴關係向下游傳播，導致多個下游服務同時出現異常徵兆，此「連鎖效應」往往掩蔽真實根因 \cite{train_ticket}。根因分析領域長期面臨\textbf{「表達力（Expressiveness）與可解釋性（Interpretability）的兩難」}：\textbf{統計與關聯分析方法}（如 BARO \cite{baro2023}）具可解釋性，但缺乏擬合複雜非線性傳播的能力；\textbf{深度學習與 GNN 方法}具強大擬合力，但其黑箱特質使維運人員無法信任，且缺乏領域歸納偏差致使推廣能力受限。特別是在\textbf{高風險（High-stakes）場景}中，即系統故障可能導致重大經濟損失、服務中斷或安全風險的關鍵維運環境（如金融交易系統、醫療資訊系統、電商核心服務），模型決策的可解釋性與可信任性成為部署的必要條件，而非僅為加分項。

傳統 GNN 依賴黑箱 MLP（Multi-Layer Perceptron）進行特徵轉換，致使維運人員難以信任其判定路徑，無法驗證模型決策是否遵循合理之因果傳播。更關鍵的是，MLP 常用的 ReLU 等啟動函數在模擬微服務異常傳播中特定的物理特性（如平滑飽和效應、閾值觸發機制）時可能存在局限。相較之下，KAN 的加法結構與一維函數特性提供視覺化驗證可能，使維運工程師能直接檢視異常傳播的函數形狀。更重要的是，KAN 透過 B-Spline 或 Chebyshev 多項式等基底函數，具備\textbf{平滑性（Smoothness）}和\textbf{局部性（Locality）}的\textbf{歸納偏差（Inductive Bias）}，使其在捕捉連續變化的異常信號時，較 MLP 更符合真實系統的物理特性。

\subsection{研究貢獻}

為打破上述兩難，本研究提出 \textbf{GNN\_KAN}——一基於 Kolmogorov–Arnold 表示定理的創新根因分析架構。GNN\_KAN 深度整合 KAN 於 GNN 消息傳遞機制中，利用 Kolmogorov–Arnold 表示定理突破表達力與可解釋性的兩難局面。本架構透過將高維函數分解為可視化的單變量函數組合，\textbf{在保留深度學習強大擬合能力的同時，賦予模型具「白箱級」的內在可解釋性（Intrinsic Interpretability）}，實現從隱式特徵到顯式函數的範式轉移。本研究主要貢獻如下：

\textbf{貢獻一：內在可解釋之圖神經網路架構}。本研究創新地將 Kolmogorov-Arnold Networks（KAN）整合至 GNN 之消息傳遞機制中，利用 Kolmogorov–Arnold 表示定理將多變量非線性函數分解為一維可學習函數之組合。相較於傳統 MLP 採用黑箱非線性啟動函數，KAN 以可參數化之一維基底函數（如 B-spline、Chebyshev 多項式、Fourier 基底、PQC\_GPU）為基礎，透過少量可學習係數實現函數組合，使模型兼具資料驅動的表達力與可視化的函數形狀。KAN 賦予模型\textbf{細粒度、逐步（step-wise）的透明度}：維運專家可在\textbf{單步消息傳遞層級}上檢查每條邊與節點更新的傳播反應曲線是否符合物理直覺（如飽和效應、門檻效應），而非聲稱整個端到端決策過程具可閉式解析解。

\textbf{貢獻二：領域感知特徵工程系統化框架}。本研究建構一套領域感知（Domain-aware）的特徵工程系統化框架，系統化地定義 RCA 場景下之特徵設計準則，專注於與根因定位高度相關且在實務應用中可穩定量測之訊號。此框架系統性規避特徵中的\textbf{後驗洩漏}風險，包括：（1）僅使用事件窗口內可觀測之前饋特徵，避免未來資訊洩漏；（2）聚焦可傳播之訊號與結構關係，排除與標註耦合之後驗指標；（3）採用一致之標準化與對齊策略，降低跨服務與跨場景之分布偏移影響。

\textbf{貢獻三：系統化基底函數選擇策略}。本研究基於廣泛的實證研究，揭示了不同基底函數（Chebyshev、B-spline、Fourier、PQC）與系統拓樸複雜度之間的內在映射關係\textbf{（如 PQC\_GPU 在超大規模系統中展現 8 倍性能優勢）}，為實務應用提供系統化的選擇準則。

\textbf{貢獻四：全面實證評估與可重現研究流程}。本研究於多個真實微服務基準（Online Boutique \cite{online_boutique}、Sock Shop \cite{sock_shop}、Train Ticket \cite{train_ticket}）及 RCAEval 平台資料集上進行全面評估，涵蓋不同拓樸結構、負載模式與異常型態。實驗結果顯示，GNN\_KAN 在 Top-k 準確度、Precision@k 等指標上展現優異性能，特別在代碼級故障場景中，實驗結果顯示皆優於其他相關方法。

本文其餘章節組織如下：第 2 節回顧相關研究；第 3 節詳述方法論；第 4 節說明實驗架構；第 5 節報告實驗結果與分析；第 6 節討論研究洞察與限制；第 7 節總結研究並提出未來研究方向。

\section{相關研究}

根因分析為分散式系統可靠性工程之核心問題，近年吸引學術界與產業界廣泛關注。本節回顧與本研究相關之方法類別，並闡述其與 GNN\_KAN 之關聯性與差異性。

多項研究利用服務依賴圖或呼叫圖建模異常傳播過程，如 \textbf{MicroRank} \cite{microrank2021} 利用隨機遊走與擴展頻譜分析定位延遲問題、\textbf{T-Rank} \cite{trank2021} 採用輕量級頻譜分析、\textbf{CloudRanger} \cite{cloudranger2018} 基於路徑分析進行呼叫鏈追蹤，以及 \textbf{BARO} \cite{baro2023} 採用貝葉斯變化點檢測方法。此外，基於 PageRank 的圖分析方法（如 PC\_PageRank）利用圖結構的中心性指標進行根因排名。此類方法具備直觀性與可解釋性，然面臨以下挑戰：（1）難以學習或初始化邊權重以精確反映真實傳播強度；（2）在多跳傳播場景中，下游節點與系統異常之統計相關性可能高於上游真實根因，導致誤判；（3）對觀測雜訊、非平穩性較為敏感，且通常假設線性或簡化之因果結構，難以捕捉複雜之非線性傳播模式。GNN\_KAN 沿襲了圖結構建模的思想，透過可學習之消息傳遞機制與 KAN 之非線性轉換，實現自適應之邊權與節點表徵調整，有效區分直接影響與間接傳播。

近年深度學習被引入維運領域，圖神經網路（GNN, Graph Neural Network）因能處理結構化資料而受到關注，已有研究將 GNN 應用於微服務 RCA \cite{gnn_rca2023}。然而，現有 GNN 方法多使用標準的非線性層（如 ReLU、GeLU）構建的 MLP，缺乏可解釋性；且在訓練資料有限或分布偏移時，推廣能力不足。實驗結果顯示，傳統 GNN (MLP) 在複雜場景中幾乎失效（Train-Ticket: P@1=0.024），證實了黑箱 MLP 在捕捉異常傳播物理特性時的局限性。此外，\textbf{Nezha} \cite{nezha2023} 與 \textbf{EADRO} \cite{eadro2023} 利用多模態資料（Metrics, Logs, Traces）進行端到端故障診斷，提升了診斷的魯棒性（Robustness）。因果推論方法如 \textbf{CausalRCA} \cite{causalrca2023} 試圖從觀測資料中重建因果圖，在可解釋性上具優勢，但在高頻、高維、非平穩的維運場景中面臨計算成本高、樣本需求大、假設難以滿足等挑戰。GNN\_KAN 的創新在於以 KAN 取代黑箱非線性，將領域知識以基底函數形狀嵌入模型，兼顧表達力與可解釋性，並提供一種在結構不確定下的近似推理框架。

知識引導學習（如 Physics-Informed Neural Networks、符號迴歸）試圖將領域知識注入模型，以提升樣本效率與推廣能力。Kolmogorov-Arnold Networks（KAN）\cite{kan2024} 是近期基於 Kolmogorov–Arnold 表示定理 \cite{kolmogorov1957} 提出的神經網路架構，將多變量函數分解為可學習的一維函數組合，使非線性轉換既具表達力又可解釋。GNN\_KAN 將 KAN 引入圖神經網路的消息傳遞中，旨在填補內在可解釋 GNN 應用於 RCA 領域的研究缺口。

近期研究如 \textbf{RCACopilot} \cite{rcacopilot2024} 與 \textbf{Xpert} \cite{xpert2024} 開始探索利用大型語言模型（LLM, Large Language Model）進行故障診斷與事件管理，透過自然語言處理提供根因分析建議。然而，LLM 在處理高頻數值型時序資料時仍面臨推理成本高、幻覺問題與可解釋性不足等挑戰，且上下文長度與數值精度限制使其難以直接對長時間窗的高維遙測序列進行精確推理，這使得基於數值計算與可視化函數的 GNN\_KAN 在即時性、精確性與可信任性上仍具明顯優勢，更適合作為數值型 RCA 引擎，而將 LLM 定位為輔助解釋與互動界面。

\section{方法論}

本節首先嚴謹定義根因分析問題，明確輸入、輸出與主要挑戰；隨後闡述 GNN\_KAN 之整體架構及各模組如何協同解決 RCA 問題；最後詳述圖建模策略、KAN 模組設計、RCA-aware 特徵工程與學習目標。

\subsection{問題定義與挑戰}

\textbf{根因分析（Root Cause Analysis, RCA）}旨在當分散式系統發生異常事件時，從大量可能之故障點中識別最可能觸發該異常之根本原因。形式化而言，給定以下輸入：

\begin{itemize}
\item \textbf{節點集合} $V = \{v_1, v_2, ..., v_n\}$，對應系統中的服務或元件（如微服務、資料庫、訊息佇列等）；
\item \textbf{有向邊集合} $E \subseteq V \times V$，表示服務間的依賴與呼叫關係，$(u, v) \in E$ 代表 $u$ 呼叫或影響 $v$；
\item \textbf{時間窗口} $W = [t_{\text{start}}, t_{\text{end}}]$，涵蓋異常事件的發生與觀測期間；
\item \textbf{節點觀測} $X = \{x_v(t) | v \in V, t \in W\}$，其中 $x_v(t) \in \mathbb{R}^d$ 為節點 $v$ 在時間 $t$ 的 $d$ 維觀測向量，包含指標（metrics，如延遲、錯誤率、吞吐量、CPU/記憶體使用）、日誌統計（logs，如異常日誌數、特定錯誤碼出現次數）、追蹤資訊（traces，如請求數、呼叫延遲分布）等；
\item \textbf{異常事件標識} $A$，指示系統在窗口 $W$ 內發生異常（如 SLA（Service Level Agreement，服務水準協議）違反、關鍵交易失敗、整體延遲飆升）。
\end{itemize}

\textbf{研究目標}為學習一排名函數 $f: (V, E, X) \to \mathbb{R}^{|V|}$，為每個節點 $v$ 計算根因分數 $s_v$，並輸出按 $s_v$ 降序排列之節點清單 $R$。理想情況下，真實根因節點 $v^*$ 應出現在 $R$ 之前列（Top-k，其中 $k << |V|$），使維運工程師能快速定位並修復問題。

\textbf{核心技術挑戰}包括：（1）\textbf{異常傳播之連鎖掩蔽效應}：根因節點之異常會沿依賴鏈向下游傳播，導致多個下游節點同時呈現異常徵兆，下游節點之異常徵兆可能更為明顯，然真實根因往往位於上游；（2）\textbf{高維度觀測與雜訊干擾}：每個節點產生多維度觀測（數十至上百維），且觀測包含雜訊、尖峰干擾、週期性波動及負載相關之正常變異；（3）\textbf{動態拓樸與結構不確定性}：服務依賴圖並非靜態，可能隨版本更新、流量路由調整、動態擴縮而變化；（4）\textbf{可解釋性與實務信任度}：維運工程師需理解「為何此節點被判定為根因」，包括支持證據、傳播路徑及與其他節點之對比。

\subsection{GNN\_KAN 解決方案架構}

針對上述挑戰，GNN\_KAN 提出整合式解決方案，透過圖神經網路與 Kolmogorov–Arnold 表示網路之協同設計實現（架構如圖 \ref{fig:architecture} 所示）：

\begin{figure*}[t!]
    \centering
    \includegraphics[width=0.95\textwidth]{figure1_structure.jpg}
    \caption{\textbf{GNN\_KAN 架構概覽。} 左側為基於 RCA-aware 特徵構建的雙圖結構（相似性圖 $G_{\text{sim}}$ 與傳播圖 $G_{\text{prop}}$）；中間層展示基於 Kolmogorov–Arnold 表示的訊息傳遞機制，以可學習的一維函數（波浪線符號）取代傳統 MLP，包含 $\text{KAN}_{\text{edge}}$、$\text{KAN}_{\text{message}}$ 與 $\text{KAN}_{\text{node}}$ 三個核心模組；右側為 $\text{KAN}_{\text{fusion}}$ 融合層與讀出層，實現全鏈路白箱化的根因定位。圖中標示了兩條主要路徑：相似性路徑（處理相似性特徵 $h^{sim}$）與傳播路徑（處理傳播特徵 $h^{prop}$），最終透過融合層產生根因分數 $s_v$。}
    \label{fig:architecture}
\end{figure*}

\textbf{階段一：基於圖消息傳遞之多跳異常傳播建模}  
本研究將微服務系統表示為動態加權有向圖 $G = (V, E, w)$，其中 $w: E \to \mathbb{R}$ 為邊權函數，反映異常傳播強度。透過多層圖神經網路之消息傳遞機制（message passing），節點表徵 $h_v$ 不僅整合自身觀測 $x_v$，亦聚合來自鄰居節點之訊息。此機制使模型能捕捉多跳傳播特性：上游節點之表徵經多層聚合後影響下游節點，而學習所得之邊權與聚合權重可反映傳播路徑之重要性。相較於單節點分析或成對關聯分析，GNN 之結構化推理能有效區分直接影響與間接傳播，緩解下游掩蔽效應。

\textbf{階段二：基於 Kolmogorov–Arnold 表示之可解釋非線性建模}  
針對傳統 GNN 中 MLP 缺乏領域語義與可解釋性的問題，本研究以 Kolmogorov-Arnold Networks（KAN）取而代之。KAN 基於 Kolmogorov–Arnold 表示定理，將多變量函數分解為一維可學習函數之組合。具體而言，KAN 以可參數化之一維基底函數（如 B-spline、Chebyshev 多項式、Fourier 基底、PQC\_GPU）建構函數空間，透過學習組合係數實現函數逼近。此設計帶來雙重優勢：（1）\textbf{函數分解之可解釋性}——每個一維函數之形狀可直接視覺化與驗證，符合「異常越強、傳播越顯著」等領域直觀；（2）\textbf{降低參數複雜度與過擬合風險}——相較於高維 MLP，一維函數組合於參數效率與推廣能力上更優。

\textbf{階段三：RCA-aware 特徵工程框架}  
特徵設計是 RCA 成敗的關鍵因素。本研究提出「RCA-aware」設計準則：僅使用與根因定位直接相關且於實務應用中可穩定量測之特徵。具體準則包括：（1）\textbf{前饋可觀測性}——僅使用事件窗口內可前饋觀測之統計量（如延遲分位數、錯誤率、資源使用率），避免引入未來資訊或事後標註指標；（2）\textbf{傳播導向性}——聚焦可傳播之訊號（如上游延遲對下游之影響）與結構關係（依賴強度、呼叫頻率），排除與標註耦合或與根因無直接關聯之特徵；（3）\textbf{一致性標準化}——採用統一之標準化與對齊策略（如 z-score 標準化、時間對齊），降低跨服務偏差與分布偏移。此框架確保模型學習通用異常傳播模式，而非訓練集特定之偽相關。

\textbf{整體架構概述}：給定異常事件窗口，GNN\_KAN 首先將各節點觀測編碼為初始表徵 $h_v^0$；隨後透過 $L$ 層消息傳遞（$L=3$ 為典型配置），每層中邊上訊息經 KAN\_edge 轉換，節點聚合鄰居訊息後，先以 KAN\_node 轉換節點表徵，再以 KAN\_message 處理聚合訊息並透過殘差連接整合；最終讀出層（readout）將 $h_v^L$ 映射為根因分數 $s_v$，並按降序輸出排名清單。訓練階段以已標註根因窗口為監督信號，最小化排名損失與正規化項；推論階段對新異常事件執行前向傳播並輸出候選清單。

\subsection{圖建模與消息傳遞機制}

GNN\_KAN 之整體架構如圖 1 所示（見「GNN\_KAN 整體架構圖」之圖片），包含三個核心組件：圖建模與消息傳遞機制、Kolmogorov-Arnold Networks 模組、以及 RCA-aware 特徵工程框架。

\subsubsection{圖表示與初始化}

本研究將分散式系統於異常事件窗口內表示為動態加權有向圖 $G_t = (V, E, W_t)$，其中節點集合 $V = \{v_1, ..., v_n\}$ 對應各服務或元件，邊集 $E \subseteq V \times V$ 由服務依賴關係與呼叫追蹤資訊萃取而得。邊權矩陣 $W_t \in \mathbb{R}^{|V| \times |V|}$ 初始化時基於服務間呼叫頻率進行歸一化（Normalization，將呼叫次數除以最大呼叫次數，確保權重落在 $[0, 1]$ 區間），訓練過程中透過可學習參數自適應調整，以反映異常傳播之實際強度。此初始化策略確保模型從合理的結構先驗開始學習，加速收斂並提升穩健性。

\textbf{初始節點編碼}：對於節點 $v_i \in V$，本研究將其於觀測窗口內之時序資料透過特徵萃取程序轉換為固定維度特徵向量 $\bar{x}_{v_i} \in \mathbb{R}^{d_{\text{feat}}}$，再經線性變換與 LayerNorm 得到初始表徵 $h_{v_i}^{(0)} \in \mathbb{R}^{d_h}$，其中 $d_h$ 為隱藏表徵維度（典型值為 64）。

\textbf{多層消息傳遞機制}：GNN\_KAN 透過 $L$ 層消息傳遞機制捕捉多跳異常傳播特性。於第 $l$ 層（$l = 1, ..., L$），依序執行以下操作：

\textbf{步驟一：邊訊息生成}。對於邊 $(v_u \to v_v) \in E$，本研究將源節點表徵 $h_u^{(l-1)}$ 與邊特徵 $e_{uv}$（包含邊權與呼叫頻率等資訊）串接後，經 KAN 模組 $\text{KAN}_{\text{edge}}$ 轉換為訊息 $m_{u \to v}^{(l)}$，此過程利用可解釋的一維函數組合來建模異常傳播的非線性模式。

\textbf{步驟二：節點訊息聚合}。節點 $v$ 接收所有入鄰居 $\mathcal{N}_{\text{in}}(v)$ 之訊息，透過注意力機制進行加權聚合。注意力權重 $\alpha_{uv}^{(l)}$ 根據源節點與目標節點之表徵相似度計算，使模型自動識別對節點 $v$ 影響最顯著之鄰居，反映異常傳播路徑之重要性。聚合後之訊息記為 $\tilde{m}_v^{(l)}$。

\textbf{步驟三：節點表徵更新}。節點更新程序分為兩階段：（1）以 $\text{KAN}_{\text{node}}$ 對節點表徵 $h_v^{(l-1)}$ 進行非線性轉換，得到 $\tilde{h}_v^{(l)}$；（2）以 $\text{KAN}_{\text{message}}$ 處理聚合訊息 $\tilde{m}_v^{(l)}$，並透過殘差連接（殘差權重 $\alpha=0.3$）與轉換後之表徵整合，得到更新後之表徵 $h_v^{(l)}$。此設計使模型先對節點表徵進行非線性轉換，再整合來自鄰居之訊息，最後透過殘差連接穩定訓練過程並促進梯度流動。

透過 $L$ 層疊加（依據實驗結果，$L=3$ 為典型配置），每個節點之表徵 $h_v^{(L)}$ 融合自身觀測與 $L$ 跳鄰居之訊息，使模型能有效捕捉深層依賴關係與異常傳播路徑。

\textbf{讀出層與根因評分}：最終層表徵 $h_v^{(L)}$ 經讀出層（線性變換 + sigmoid 函數）映射為根因分數 $s_v \in (0, 1)$。高分數表示節點 $v$ 更可能為根因。所有節點按 $s_v$ 降序排列，輸出候選清單 $R$。

\textbf{雙圖構建策略}：本研究採用雙圖策略以整合結構相似性與異常傳播資訊。\textbf{相似性圖 $G_{\text{sim}}$} 基於特徵餘弦相似度之無向圖，邊權 $w_{\text{sim}}(u, v)$ 定義為兩節點特徵向量之餘弦相似度減去閾值 $\tau_{\text{sim}}$（預設 0.5）後取最大值，捕捉結構相似性。\textbf{傳播圖 $G_{\text{prop}}$} 基於時滯相關性之有向圖，捕捉異常傳播方向。對於節點對 $(u, v)$，本研究計算其異常信號之最大時滯相關性 $\rho_{\text{lag}}(u, v)$。若相關性超過閾值 $\tau_{\text{prop}}$ 且 $u$ 之異常發生時間早於 $v$，則建立有向邊 $u \to v$，權重 $w_{\text{prop}}(u, v)$ 與時滯相關性成正比，與時間差成反比，編碼時序一致性。於前向傳播過程中，本研究對兩圖分別執行消息傳遞，得到基於相似性的表徵 $h_v^{\text{sim}}$ 與基於傳播的表徵 $h_v^{\text{prop}}$。為維持架構的一致性與完全可解釋性，本研究採用 \textbf{KAN 融合層（KAN Fusion Layer）} 進行特徵整合：

\begin{equation}
h_v^{\text{final}} = \text{KAN}_{\text{fuse}}\left( h_v^{\text{sim}} \oplus h_v^{\text{prop}} \right)
\end{equation}
（基於 KAN 架構 \cite{kan2024} 的融合層設計）

其中 $\oplus$ 代表\textbf{特徵維度上的張量串接（Concatenation along feature dimension）}，即 $[h_v^{\text{sim}}; h_v^{\text{prop}}] \in \mathbb{R}^{2d}$。透過 KAN 的可加性結構，模型能以非線性且可視化的方式，動態權衡結構相似性與異常傳播資訊的重要性。此設計消除了模型中最後的黑箱組件，確保從特徵輸入、訊息傳遞到最終融合的每一環節，皆具備數學上的內在可解釋性。

\textbf{KAN 與雙圖策略的內在連結}：值得注意的是，KAN 的一維函數分解特性使其在處理異質圖結構時具有獨特的適應性。在\textbf{傳播圖}中，KAN 傾向於學習具有「門檻觸發（Threshold Triggering）」的函數形狀（如 Sigmoid 狀），以模擬異常的觸發機制；而在\textbf{相似性圖}中，KAN 則傾向於學習「平滑映射（Smooth Mapping）」，以捕捉功能相似服務間的共變異性。這種\textbf{結構依賴的函數自適應性（Structure-dependent Function Adaptivity）}是傳統共享權重的 MLP 所難以實現的，因為 MLP 的權重在兩個圖結構間是共享的，無法根據圖的語義特性自適應調整函數形狀。

\subsection{RCA-aware 特徵工程框架}

特徵設計是決定 RCA 模型性能的關鍵因素，不當之特徵會引入資料洩漏（data leakage）或分布偏移（distribution shift），導致訓練集高準確度然部署時失效。本研究提出 \textbf{RCA-aware} 特徵設計原則，並基於實作架構提供系統化定義。

\subsubsection{RCA-aware 特徵設計準則與系統化定義}

本研究遵循四大 RCA 特徵設計準則，系統性地消除後驗洩漏風險：（1）僅使用前饋特徵（事件窗口 $W$ 及歷史回溯窗口 $W_{\text{hist}}$ 內可得，嚴格排除未來資訊），滿足因果約束；（2）聚焦可傳播訊號（異常強度、時序變化、傳播關聯性）；（3）排除與標註耦合的後驗指標；（4）實施一致的 z-score 標準化策略。對每個服務節點 $v$ 萃取 48 維特徵向量 $\bar{x}_v$，包含：SLI（Service Level Indicator，服務水準指標）關聯特徵（同步/時滯相關性、滯後量）、時序差異特徵（故障前後統計量變化）、異常偵測特徵（峰值/持續異常強度）與故障類型特化特徵（延遲分位數、變異係數、偏態等），並進行標準化處理與 NaN/Inf 值處理以確保數值穩定性。

\subsection{Kolmogorov-Arnold Networks 模組設計}

傳統神經網路使用固定形狀之啟動函數（如 ReLU、Sigmoid、Tanh）與多層感知器（MLP），此類黑箱非線性模組雖計算簡單且表達力強，然缺乏可解釋性且易過度擬合。\textbf{Kolmogorov-Arnold Networks（KAN）}基於 Kolmogorov–Arnold 表示定理，提供可學習且可解釋之函數逼近替代方案。

\subsubsection{Kolmogorov–Arnold 表示定理與 RCA 的內在連結}

\textbf{Kolmogorov–Arnold 表示定理}（Kolmogorov, 1957; Arnold, 1957）指出：任意連續多變量函數可分解為一維連續函數之組合 $f(x_1, \ldots, x_n) = \sum_{q=0}^{2n} \phi_q\left( \sum_{p=1}^{n} \psi_{p,q}(x_p) \right)$。KAN 將此數學原理應用於神經網路：以可參數化之一維基底函數取代固定之非線性啟動（如 ReLU），透過學習函數組合實現多變量映射。微服務異常傳播遵循特定物理規律（飽和效應、閾值效應、平滑漸變），MLP 常用的 ReLU 等啟動函數在模擬這些平滑特性時可能存在局限。KAN 透過 B-Spline 或 Chebyshev 等基底函數，具備\textbf{平滑性與局部性的歸納偏差}，更符合系統物理直覺。KAN 應用於 RCA 具備三重優勢：（1）內在可解釋之函數分解（一維函數可視覺化，驗證單調性、門檻性、飽和性）；（2）參數效率（$O(d_{\text{in}} \times K)$ vs. $O(d_{\text{in}} \times d_{\text{out}})$，降低 $O(d_{\text{out}}/K)$ 倍，減少過度擬合）；（3）領域知識注入靈活性（基底函數選擇編碼領域知識）。

\subsubsection{一維函數參數化與數學定義}

KAN 之核心為將多變量函數分解為一維函數之組合。於實作中，本研究對每個輸入維度 $x_i$ 建立一維可學習函數 $\psi_i: \mathbb{R} \to \mathbb{R}$，最終輸出為此一維函數之線性組合。形式化而言，對於輸入 $z = [z_1, \ldots, z_{d_{\text{in}}}]^T \in \mathbb{R}^{d_{\text{in}}}$，KAN 層之輸出定義為：

\begin{equation}
\phi(z) = \sum_{i=1}^{d_{\text{in}}} \sum_{k=1}^{K} \alpha_{i,k} \cdot b_k(z_i)
\end{equation}
（基於 Kolmogorov–Arnold 表示定理 \cite{kolmogorov1957} 的 KAN 層定義 \cite{kan2024}）

其中 $K$ 為一維基底函數之數量（典型值為 5-10），$b_k: \mathbb{R} \to \mathbb{R}$ 為一維基底函數，$\alpha_{i,k} \in \mathbb{R}$ 為可學習係數。此結構遵循 Kolmogorov–Arnold 表示之精神：多變量映射 $\phi: \mathbb{R}^{d_{\text{in}}} \to \mathbb{R}$ 分解為 $d_{\text{in}}$ 個一維函數 $\psi_i(z_i) = \sum_{k=1}^{K} \alpha_{i,k} \cdot b_k(z_i)$ 之加權和。

\textbf{參數效率分析}：相較於傳統 MLP 需要 $O(d_{\text{in}} \times d_{\text{out}})$ 個參數，KAN 僅需 $O(d_{\text{in}} \times K)$ 個參數（假設 $K << d_{\text{out}}$），參數複雜度降低 $O(d_{\text{out}}/K)$ 倍。此參數效率優勢於標註稀缺之 RCA 場景中尤為重要，降低過度擬合風險並提升推廣能力。

\textbf{基底函數族選擇}：本研究支援四種基底函數，各具不同特性與「表達力–可解釋性」權衡：\textbf{（1）Chebyshev 多項式基底}（預設）：$b_k^{\text{Cheb}}(z) = T_k(\text{tanh}(z))$，具備最優逼近性質與數值穩定性，適合扇出型拓樸；\textbf{（2）B-spline 基底}：$b_k^{\text{Spline}}(z) = B_{k,p}(z)$（$p=3$ 立方樣條），具備局部支撐性，適合線性拓樸之平滑傳播；\textbf{（3）Fourier 基底}：$b_k^{\text{Fourier}}(z) = \cos(k \pi z)$ 或 $\sin(k \pi z)$，適合捕捉週期性異常模式；\textbf{（4）參數化量子電路基底（PQC\_GPU, Parameterized Quantum Circuit）}：作為\textbf{針對超大規模、深層複雜拓樸的高表達力選項}，用於處理傳統多項式基底出現飽和的極端情境。對於超大規模、深層依賴系統（如 Train Ticket），節點間高階交互作用導致特徵空間指數級膨脹。PQC 利用量子態希爾伯特空間（Hilbert Space）的張量積特性，以\textbf{多項式級}參數複雜度模擬高維交互：
\begin{equation}
b_k^{\text{PQC}}(z) = \langle 0|^{\otimes n} U^\dagger_{\text{var}}(\boldsymbol{\theta}) U_{\text{enc}}(z) \, M \, U_{\text{enc}}(z) U_{\text{var}}(\boldsymbol{\theta}) |0\rangle^{\otimes n}
\end{equation}
（基於變分量子電路（VQC, Variational Quantum Circuit）架構 \cite{pennylane} 的 PQC 基底函數定義）
其中 $U_{\text{enc}}(z)$ 為資料編碼么正算符（將輸入 $z$ 映射至量子態），$U_{\text{var}}(\boldsymbol{\theta})$ 為變分電路（可學習參數 $\boldsymbol{\theta}$），\textbf{$M$ 為測量算符}（如 Pauli-Z 張量積 $Z^{\otimes n}$）。此設計遵循變分量子電路（VQC, Variational Quantum Circuit）的標準架構：編碼、參數化變換、測量。本研究使用 PennyLane 框架實作 GPU 加速之量子電路模擬，具備 $2^Q$ 維量子態空間之指數級表示能力（$Q=4$ 量子位元），實驗證實其在超大規模複雜系統中之優勢。\textbf{需注意的是，PQC 在幾何形狀上的直觀可解釋性不如 B-spline/Chebyshev，其可解釋性主要體現在顯式張量積結構與複雜度邊界的數學分析上，因此應被視為在「表達力優先」場景下的權衡選項，而非本研究所強調之可視化函數形狀的主要來源。}

\subsubsection{KAN 層之前向傳播機制}

KAN 層之前向傳播包含三個平行計算組件：（1）\textbf{基底函數計算}：對批次輸入 $Z$ 計算所有基底函數 $b_k(Z)$，得到 $\mathcal{B}(Z) \in \mathbb{R}^{B \times d_{\text{in}} \times K}$；（2）\textbf{基底函數線性組合}：使用 Einstein 求和約定（einsum）將基底函數與可學習係數張量 $A$ 收縮，得到 $\Phi_{\text{basis}}(Z)$；（3）\textbf{可學習線性基底}：計算 $\Phi_{\text{linear}}(Z) = W_{\text{base}} \cdot Z + b_{\text{base}}$ 以保留表達力。最終輸出透過 $\tanh$ 限制範圍、縮放係數（0.5, 0.1）平衡兩組件貢獻，並以 LayerNorm 穩定訓練過程。

\textbf{於 GNN\_KAN 中之應用位置}：（1）\textbf{KAN\_edge}：於邊訊息生成階段，將 $[h_u^{(l-1)} \oplus e_{uv}]$ 經 KAN 轉換為訊息 $m_{u \to v}^{(l)}$；（2）\textbf{KAN\_node}：於節點更新階段，先對 $h_v^{(l-1)}$ 進行 KAN 轉換得到 $\tilde{h}_v^{(l)}$；（3）\textbf{KAN\_message}：處理聚合訊息 $\tilde{m}_v^{(l)}$，最後透過殘差連接整合。

\subsubsection{KAN 與 GNN 消息傳遞的理論契合度分析}

本小節從理論層面深入分析為何 KAN 的結構先驗與 GNN 消息傳遞機制在 RCA 問題上具有內在契合度，以及這種契合度如何轉化為性能優勢。

\textbf{消息傳遞的可加性結構}：在標準 GNN 的消息傳遞過程中，節點 $v$ 的表徵更新遵循以下形式：
\begin{equation}
h_v^{(l)} = \text{UPDATE}\left(h_v^{(l-1)}, \text{AGGREGATE}\left(\{m_{u \to v}^{(l)} : u \in \mathcal{N}(v)\}\right)\right)
\end{equation}
其中 $\text{AGGREGATE}$ 通常為求和（SUM）、平均（MEAN）或注意力加權求和。關鍵洞察在於：\textbf{異常傳播在微服務系統中本質上具有可加性}。當多個上游服務同時出現異常時，其對下游服務的影響近似為各上游異常影響的加權和，而非複雜的非線性交互。這與 Kolmogorov–Arnold 表示中的可加性結構（$\sum_{i=1}^{d_{\text{in}}} \psi_i(z_i)$）高度吻合。

\textbf{一維函數分解與異常傳播的單變量依賴}：微服務異常傳播的一個重要特性是：\textbf{單個上游節點的異常強度（如 CPU 使用率、延遲）對下游節點的影響往往可以建模為該異常強度的單變量函數}。例如，上游服務的延遲 $d_u$ 對下游服務延遲的影響 $f(d_u)$ 通常呈現平滑的單調遞增關係，而非依賴於多個變量的複雜交互。KAN 的一維函數分解 $\psi_i(z_i)$ 恰好捕捉了這種單變量依賴關係，而 MLP 的全連接結構則會學習不必要的多變量交互項，導致過度擬合。

\textbf{歸納偏差與物理約束的一致性}：微服務系統的異常傳播遵循特定的物理約束：（1）\textbf{平滑性}：異常強度的微小變化不應導致傳播效果的劇烈跳躍；（2）\textbf{局部性}：遠距離服務的異常對當前節點的影響應隨距離衰減；（3）\textbf{飽和性}：當異常強度超過某閾值後，進一步增加異常強度對傳播效果的邊際影響遞減。KAN 的 B-spline 和 Chebyshev 基底函數天然具備平滑性和局部性，而 MLP 的 ReLU 啟動函數則缺乏這些約束，可能學習到不符合物理直覺的階梯狀函數。

\textbf{複雜交互的處理與 PQC 的定位}：雖然大部分異常傳播遵循可加性結構，但在超大規模系統（如 Train Ticket，40+ 服務）中，確實存在高階交互作用（例如，服務 A 和服務 B 同時異常時對服務 C 的影響不等於兩者單獨影響的簡單加和）。此時，PQC\_GPU 的指數級表示能力提供了必要的建模能力。然而，\textbf{我們必須坦誠承認 PQC 的可解釋性局限}：其可解釋性主要體現在數學結構層面（張量積分解、多項式複雜度邊界），而非如 B-spline 般具備直觀的幾何形狀。因此，PQC 應被定位為一種\textbf{針對極端複雜場景的權衡方案}，在可解釋性與表達力之間做出明確取捨。在實際應用中，建議優先使用 Chebyshev 或 B-spline，僅在系統規模超過 30 個服務且出現明顯的高階交互時才考慮 PQC。

\textbf{理論優勢的實證驗證}：上述理論分析在實驗中得到了驗證。在簡單線性場景（RE1-OB）中，統計方法 BARO 因參數效率優勢而表現更優；但在複雜非線性場景（Train Ticket）中，GNN\_KAN 的優勢顯現，證實了 KAN 的歸納偏差在複雜傳播模式中的有效性。更重要的是，KAN 學習到的函數形狀（見圖 \ref{fig:kan_viz}）呈現出符合物理特性的「死區/飽和區」結構，而 MLP 則學習到無規律的雜訊，這從實證角度支持了理論分析。

\subsection{訓練目標與學習機制}

GNN\_KAN 之訓練目標整合多個子目標，兼顧準確度、可解釋性與穩健性：

\textbf{主要目標：排名損失（Ranking Loss）}  
RCA 之核心為排名而非分類，因此本研究採用排名導向之損失函數。給定一標註根因為 $v^*$ 之異常事件，本研究希望 $s_{v^*}$ 顯著高於其他節點。本研究採用 Margin Ranking Loss：

\begin{equation}
L_{\text{rank}} = \sum_{v \neq v^*} \max(0, \text{margin} - (s_{v^*} - s_v))
\end{equation}
（標準 Margin Ranking Loss 形式）

此損失函數鼓勵根因分數與非根因分數之差距大於邊際值 $\text{margin}$，對 Top-k 指標最佳化有效。

\textbf{正規化項一：邊權稀疏化（Edge Sparsity Regularization）}  
為避免模型學習過於複雜之傳播結構，本研究對邊權施加 L1 正規化：

\begin{equation}
L_{\text{edge}} = \lambda_1 \cdot \sum_{(u,v)\in E} |w_{uv}|
\end{equation}
（L1 正規化標準形式）

其中 $\lambda_1=0.01$ 為正規化係數。此正規化項鼓勵模型學習稀疏之關鍵傳播路徑，提升可解釋性並降低過度擬合風險。

\textbf{正規化項二：KAN 平滑化（KAN Smoothness Regularization）}  
為防止 KAN 基底係數過度擺動導致非線性函數不穩定，本研究對係數施加平滑約束：

\begin{equation}
L_{\text{kan}} = \lambda_2 \cdot \sum_{\text{KAN}} \sum_k |\alpha_k - \alpha_{k-1}|
\end{equation}
（平滑正規化項，用於約束 KAN 基底係數的連續性）

其中 $\lambda_2=0.001$ 為正規化係數。此正規化項保留基底函數之先驗形狀（如 B-spline 之光滑性），避免過度擬合並提升跨場景推廣能力。

\textbf{總損失函數}：$L_{\text{total}} = L_{\text{rank}} + \lambda_1 \, L_{\text{edge}} + \lambda_2 \, L_{\text{kan}}$。

\textbf{最佳化策略}：本研究使用 Adam 最佳化器，學習率採用 warm-up + cosine decay 排程以穩定訓練過程。批次化處理多個異常事件窗口，每批隨機抽樣事件並對應之圖結構。

\section{實驗架構}

本節說明以 RCAEval 平台實現之端到端實驗架構，確保方法論、資料處理與評估協定的一致性與可重現性。

\subsection{基礎資料集特性}

本研究採用 RCAEval 平台的五個核心資料集，涵蓋不同規模與故障類型。\textbf{Sock Shop 系列}（v1/v2）為線性拓樸的中等規模系統（約 10 個服務），v2 在 v1 基礎上擴展了追蹤資訊，涵蓋資源與網路故障。\textbf{RE1 系列}為標準化評估基準：RE1-OB（Online Boutique，扇出型拓樸，375 個案例）、RE1-SS（Sock Shop，線性拓樸，375 個案例）、RE1-TT（Train Ticket，40+ 服務深層複雜拓樸），涵蓋資源與網路故障。\textbf{RE2 系列}在 RE1 基礎上增加多源遙測資料（metrics+logs+traces），RE2-OB（270 個案例，77-376 個指標/服務）、RE2-SS/TT（超大規模多源資料）。\textbf{RE3 系列}專注代碼級故障（5 種類型，90 個案例），在 OB/SS/TT 三個系統上評估精細根因定位能力。此資料集選擇涵蓋從簡單到超大規模、從資源/網路到代碼級故障的完整評估光譜。

\subsection{實驗流程架構}

本研究之實驗流程遵循標準化管線，確保方法論、資料處理與評估協定之一致性與可重現性，包含以下主要階段：

\textbf{階段一：資料載入與預處理}
\begin{itemize}
\item \textbf{資料載入}：由 RCAEval 平台之 Datasets/Registry 管理，明確標示資料集版本、切分比例（訓練/驗證/測試：60\%/20\%/20\%）與來源，確保事件間之時間獨立性以避免資料洩漏。
\item \textbf{資料驗證}：檢查資料完整性、時間序列連續性與標註一致性。
\end{itemize}

\textbf{階段二：特徵工程與圖構建}
\begin{itemize}
\item \textbf{特徵處理}：落實第 3.4 節之 RCA-aware 原則（前饋可觀測、傳播導向、無標註耦合、一致正規化），對每個服務節點提取 48 維特徵向量。
\item \textbf{圖構建}：實作雙圖策略（相似性圖與傳播圖）與時滯相關性計算（第 3.3 節），並輸出標準化之圖張量（節點特徵矩陣、邊索引、邊權重）。
\end{itemize}

\textbf{階段三：模型訓練}
\begin{itemize}
\item \textbf{模型初始化}：根據設定初始化 GNN\_KAN 模型，包含多層消息傳遞（L=3）與 KAN 模組，KAN 基底函數類型（Chebyshev、B-spline、Fourier、PQC\_GPU）由設定指定。
\item \textbf{訓練執行}：採用排名損失（Margin Ranking Loss）與正規化項（邊稀疏正規化 $\lambda_1=0.01$、KAN 平滑正規化 $\lambda_2=0.001$），使用 Adam 最佳化器與 warm-up + cosine decay 學習率排程。
\item \textbf{驗證與早停}：於驗證集上監控性能，實施早停機制避免過度擬合。
\end{itemize}

\textbf{階段四：評估與分析}
\begin{itemize}
\item \textbf{推論執行}：對測試集執行前向傳播，計算各節點之根因分數並生成排名清單。
\item \textbf{指標計算}：計算 Top-k Accuracy、Precision@k、Recall@k、mAP、MRR 等評估指標。
\item \textbf{結果輸出}：生成結構化產物，包含模型權重、評估指標（CSV/JSON）、Top-k 排名清單、KAN 函數曲線可視化、邊權熱圖與完整日誌。
\end{itemize}

\subsection{模型與學習設定}

\textbf{模型架構}：GNN\_KAN 包含多層消息傳遞（L=3）與 KAN 模組（KAN\_edge、KAN\_node、KAN\_message），KAN 基底函數類型（Chebyshev、B-spline、Fourier、PQC\_GPU）與係數之設置受設定控制。

\textbf{訓練機制}：採用排名損失（Margin Ranking Loss），並加入邊稀疏正規化（$\lambda_1=0.01$）與 KAN 平滑正規化（$\lambda_2=0.001$）。最佳化器使用 Adam，學習率採用 warm-up + cosine decay 排程，批次大小為 32。

\textbf{隨機性控制}：透過平台統一設定隨機種子，涵蓋資料抽樣、權重初始化與後端運算，確保實驗可重現性。

\subsection{實驗設定與執行}

\textbf{統一協定}：所有方法使用相同之資料切分、窗口長度（異常觸發前後各 5 分鐘）、特徵萃取流程。圖結構由服務依賴與調用追蹤資訊萃取，邊權初始化為正規化之呼叫頻率。\textbf{基線方法設定}：GNN (MLP) 採用與 GNN\_KAN 相同的架構（3 層消息傳遞、表徵維度 64），僅將 KAN 模組替換為標準 MLP（ReLU 啟動函數、隱藏層維度 64），確保公平對比。PC\_PageRank 基於服務依賴圖計算 PageRank 分數，並結合異常強度進行加權排名。

\textbf{因果基線說明}：雖然相關研究中提及 CausalRCA 等因果推論方法，但在本研究所採用之大規模、多源遙測資料場景下，其計算複雜度與樣本需求遠高於可接受範圍。初步嘗試顯示，直接在 RE2/RE3 級別資料集上套用完整因果結構學習流程在時間與記憶體成本上均難以承受，且需要大量額外先驗假設。更重要的是，多數因果發現演算法假設變數關係構成有向無環圖（DAG），而真實微服務依賴圖普遍存在循環呼叫與重試機制，違反其基礎假設，導致直接套用經典因果演算法的理論基礎不足。基於公平性與可重現性考量，本研究聚焦於在相同特徵與圖結構設定下，系統性比較統計方法（BARO）、傳統 GNN (MLP)、圖分析方法（PC\_PageRank）與所提 GNN\_KAN，而將因果方法留作未來工作的一部分。

\textbf{超參數}：GNN\_KAN 之超參數（層數 L=3、表徵維度 d=64、KAN 基底數 K=8、學習率 0.001、正規化係數 $\lambda_1=0.01, \lambda_2=0.001$）透過驗證集調整。

\textbf{計算環境}：實驗於配備 NVIDIA 2080ti GPU 之伺服器上執行，訓練與推論時間均記錄以評估效率。

\textbf{可重現性（RCAEval 平台）}：為使上述方法能於不同資料集與場景中以一致方式執行與比較，本研究將整體流程編排於 RCAEval 平台，包含標準化之資料載入（Datasets/Registry）、RCA-aware 特徵處理（Features/Processors）、雙圖構建（Graphs/Builders）與模型訓練評估（Runners/Evaluators）。所有實驗透過此平台執行，採用統一之設定管理、資料載入與版本控管、隨機性控制、實驗執行器與產物追蹤協定。

\section{實驗結果}

本節報告 GNN\_KAN 與多個基線方法之實驗結果，涵蓋五個核心資料集（sock-shop-2、RE1、RE2、train-ticket）之全面評估。本研究比較 GNN\_KAN 於不同基底函數（Chebyshev、B-spline、Fourier、PQC\_GPU）下之表現，並與以下基線方法進行對比：（1）\textbf{BARO}——基於貝葉斯變化點檢測的統計方法；（2）\textbf{GNN (MLP)}——使用傳統 MLP 的 GNN 基線，採用與 GNN\_KAN 相同的 RCA-aware 特徵提取與圖結構，僅將 KAN 模組替換為標準 MLP，以公平評估 KAN 的相對優勢；（3）\textbf{PC\_PageRank}——基於 PageRank 的圖分析方法，利用服務依賴圖進行根因排名。此對比設計旨在深入分析各方法於不同故障類型與系統規模下之優劣，並驗證 KAN 相較 MLP 的架構優勢。

\begin{table*}[htbp]
\centering
\caption{GNN\_KAN 與基線方法在不同複雜度場景下的性能比較：GNN (MLP) 與 PC\_PageRank 在複雜場景中的性能崩潰，凸顯 KAN 架構與可學習消息傳遞之必要性}
\label{tab:overall_results}
\small
\begin{tabular}{ll|cccc|cccc|cccc|cccc}
\toprule
\multirow{2}{*}{\textbf{場景特徵}} & \multirow{2}{*}{\textbf{資料集}} & \multicolumn{4}{c|}{\textbf{GNN\_KAN}} & \multicolumn{4}{c|}{\textbf{BARO}} & \multicolumn{4}{c|}{\textbf{GNN (MLP)}} & \multicolumn{4}{c}{\textbf{PC\_PageRank}} \\
\cmidrule(lr){3-6} \cmidrule(lr){7-10} \cmidrule(lr){11-14} \cmidrule(lr){15-18}
& & \textbf{P@1} & \textbf{P@3} & \textbf{P@5} & \textbf{Avg@5} & \textbf{P@1} & \textbf{P@3} & \textbf{P@5} & \textbf{Avg@5} & \textbf{P@1} & \textbf{P@3} & \textbf{P@5} & \textbf{Avg@5} & \textbf{P@1} & \textbf{P@3} & \textbf{P@5} & \textbf{Avg@5} \\
\midrule
\multirow{2}{*}{\makecell[l]{\textbf{簡單/線性場景}\\\scriptsize (資源故障，\\線性依賴)}} 
& \textbf{RE1-OB} & 0.432 & 0.624 & 0.632 & 0.584 & \cellcolor{gray!20}\textbf{0.736} & \cellcolor{gray!20}\textbf{0.928} & \cellcolor{gray!20}\textbf{0.984} & \cellcolor{gray!20}\textbf{0.896} & 0.088 & 0.352 & 0.584 & 0.349 & 0.080 & 0.216 & 0.448 & 0.242 \\
& \textbf{RE1-SS} & 0.264 & 0.720 & 0.952 & 0.680 & 0.280 & 0.728 & 0.920 & 0.664 & 0.056 & 0.280 & 0.600 & 0.298 & - & - & - & - \\
\midrule
\multirow{2}{*}{\makecell[l]{\textbf{複雜/非線性場景}\\\scriptsize (大規模拓樸，\\多跳傳播)}} 
& \textbf{Train-Ticket} & \cellcolor{gray!20}\textbf{0.256} & \cellcolor{gray!20}\textbf{0.376} & \cellcolor{gray!20}\textbf{0.456} & \cellcolor{gray!20}\textbf{0.365} & 0.136 & 0.280 & 0.392 & 0.272 & 0.024 & 0.048 & 0.080 & 0.051 & 0.110 & 0.140 & 0.152 & 0.128 \\
& \textbf{Sock-Shop-2} & 0.272 & 0.784 & 0.936 & 0.678 & 0.280 & 0.728 & 0.920 & 0.664 & 0.112 & 0.328 & 0.624 & 0.344 & 0.030 & 0.350 & 0.490 & 0.344 \\
\midrule
\multirow{2}{*}{\makecell[l]{\textbf{高維/多源雜訊}\\\scriptsize (Metrics+Logs\\+Traces)}} 
& \textbf{RE2-OB} & \cellcolor{gray!20}\textbf{0.478} & 0.811 & 0.856 & \cellcolor{gray!20}\textbf{0.758} & 0.144 & \cellcolor{gray!20}\textbf{0.878} & \cellcolor{gray!20}\textbf{0.944} & 0.742 & 0.322 & 0.789 & 0.844 & 0.696 & 0.080 & 0.484 & 0.628 & 0.492 \\
& \textbf{RE2-SS} & \cellcolor{gray!20}\textbf{0.333} & 0.811 & 0.900 & 0.693 & 0.067 & \cellcolor{gray!20}\textbf{0.956} & \cellcolor{gray!20}\textbf{0.978} & \cellcolor{gray!20}\textbf{0.764} & 0.156 & 0.589 & 0.878 & 0.560 & 0.150 & 0.425 & 0.667 & 0.414 \\
\bottomrule
\end{tabular}
\vspace{0.2cm}

\small \textit{註：灰底標示該場景類別的最佳方法。GNN\_KAN 使用 Chebyshev（除 Train-Ticket 用 PQC\_GPU）。GNN (MLP) 為使用傳統 MLP 的 GNN 基線，採用相同 RCA-aware 特徵提取。PC\_PageRank 為基於 PageRank 的圖分析方法。P@1 是最關鍵指標（前 1 名準確度）。}
\end{table*}

\subsection{消融研究與方法論邊界分析}

實驗結果（表 \ref{tab:overall_results}）從「黑箱 MLP」、「傳統圖算法」與「統計方法」三個角度，系統性地揭示了各類方法在 RCA 任務中的\textbf{適用邊界與失效模式}，並回答了「為什麼需要 KAN 而不是單純的 MLP 或 PageRank」這個關鍵問題。除與統計基線 BARO 的比較外，本研究特別納入兩個強對照組：\textbf{GNN (MLP)}——僅將 GNN\_KAN 中的 KAN 模組替換為標準 MLP 的消融版本，以及 \textbf{PC\_PageRank}——基於 PageRank 的圖分析方法，用以檢驗「僅依賴靜態圖拓樸」是否足以支撐 RCA。

\textbf{一、黑箱 MLP 的崩潰：從「略遜一籌」到「幾乎失效」}

GNN (MLP) 直接檢驗了「若僅替換非線性模組為 MLP，是否能取得類似性能」的假設。在看似簡單的 RE1-OB 場景中，GNN (MLP) 的 P@1 僅為 0.088，遠低於 GNN\_KAN 的 0.432，顯示即便在結構與特徵相同的情況下，黑箱 MLP 仍難以穩定學習異常傳播模式。更為關鍵的是，在複雜拓樸場景 RE1-TT 中，GNN (MLP) 的 P@1 下降至 \textbf{0.024}，幾乎接近隨機猜測水準，而 GNN\_KAN 仍維持在 0.224。這種「性能崩潰」現象說明，標準 MLP 所採用的 ReLU/GeLU 等全局啟動函數，在面對多跳傳播、弱信號疊加與不平衡樣本時，容易出現梯度消失或嚴重過擬合；相較之下，KAN 透過一維基底函數與可視化的傳播反應曲線，為模型注入平滑性與局部性之歸納偏差，使其在深層依賴圖上仍能「學得動」且具備穩定可用性。

\textbf{二、靜態圖分析的失效：僅憑拓樸結構遠遠不夠}

PC\_PageRank 回答了另一個常見質疑：「既然你強調圖結構這麼重要，為什麼不用 PageRank/Random Walk 這類經典圖算法？」實驗結果顯示，PC\_PageRank 在 Train-Ticket 的 P@1 僅有 \textbf{0.110}，在 RE2-OB 也只有 0.080，明顯落後於 GNN\_KAN（Train-Ticket: 0.256, RE2-OB: 0.478）。這表明\textbf{僅依賴靜態拓樸與固定轉移機率，無法刻畫微服務故障的動態傳播機制}：服務間的依賴關係具有條件觸發性（僅在特定負載或故障模式下才強烈傳播），而 PageRank 假設的則是與內容無關、時間不變的轉移機率。相比之下，GNN\_KAN 透過可學習的消息傳遞機制與邊權調整，能根據異常強度與時序特徵自適應地重塑有效傳播路徑。即便在採用 PQC\_GPU 基底的設定下，模型仍然保留了\textbf{結構層級的可解釋性（Structural Interpretability）}：維運工程師可以透過檢視學習到的邊權重與注意力分佈，追蹤由根因節點出發的高權重傳播路徑，理解「異常如何在圖上擴散」，雖然對單一非線性函數形狀的幾何直觀性有所折衷。

\textbf{三、簡單場景下的誠實邊界：統計方法仍具優勢}

在結構簡單、故障模式接近線性的 RE1-OB/Online-Boutique 等場景中，BARO 仍展現出統計方法的優勢（P@1=0.736 vs GNN\_KAN 0.432–0.472）。這一結果符合奧卡姆剃刀原理：在訊號清晰且近似線性的設定下，參數更少且假設明確的統計模型更難被超越。然而，值得注意的是，即便在這類「對 BARO 友善」的場景中，GNN (MLP) 的 P@1 僅為 0.088，明顯落後於 GNN\_KAN（0.432），顯示\textbf{MLP 不僅在複雜場景中幾乎失效，在簡單場景中也難以充分利用 RCA-aware 特徵}。因此，GNN\_KAN 與 BARO 構成的是一組互補方法：前者對應高複雜度與非線性場景，後者則適用於低維線性問題，而黑箱 MLP 在兩端皆缺乏明確優勢。

\textbf{核心結論：從「可解釋性加分」到「可用性門檻」}

綜上所述，本實驗不僅劃定了統計方法、傳統 GNN (MLP)、圖算法與 GNN\_KAN 的\textbf{適用邊界}，更揭示了 KAN 在 RCA 任務中的角色：\textbf{其貢獻已不僅是提供可視化函數形狀的「解釋性加分」，而是確保模型在複雜場景下仍具備基本可用性的結構性前提}。在簡單問題的邊界內，統計方法仍是強而有力的基線；一旦跨入高維、多源、深層拓樸與代碼級故障等困難區域，GNN\_KAN 相較 GNN (MLP) 與 PC\_PageRank 展現出明顯且穩定的優勢。隨著雲原生系統持續朝向更大規模與更高複雜度演進，引入 Kolmogorov–Arnold 表示與 KAN 結構，已成為在 RCA 場景中兼顧表達力、可解釋性與可用性的\textbf{戰略性必要選擇}。

\subsection{效率與開銷：實務應用可行性}

在 NVIDIA 2080ti GPU 上，GNN\_KAN 推論時間約 50-100ms（vs. MLP 的 30-50ms），考慮 RCA 場景通常是分鐘級窗口，此延遲可接受。訓練時間增加 20-30\%，但 KAN 參數量僅為 MLP 的 60-70\%，在記憶體受限環境更具優勢。PQC\_GPU 計算開銷最高（約 Chebyshev 的 2-3 倍），但其在超大規模複雜系統的性能優勢證實此開銷合理。

\subsection{基底函數比較與選擇策略}

\begin{table*}[htbp]
\centering
\caption{不同基底函數在代表性資料集上的性能比較（Precision@1）}
\label{tab:basis_comparison}
\small
\begin{tabular}{l|cccc}
\toprule
\textbf{資料集} & \textbf{Chebyshev} & \textbf{B-spline} & \textbf{Fourier} & \textbf{PQC\_GPU} \\
\midrule
\textbf{RE1-OB} & \textbf{0.432} & 0.392 & 0.432 & 0.360 \\
\textbf{RE1-SS} & 0.264 & \textbf{0.272} & 0.224 & 0.136 \\
\textbf{RE1-TT} & 0.080 & 0.008 & 0.008 & \textbf{0.224} \\
\midrule
\textbf{RE2-OB} & 0.478 & 0.367 & \textbf{0.478} & 0.478 \\
\textbf{RE2-SS} & \textbf{0.333} & 0.289 & 0.322 & 0.200 \\
\midrule
\textbf{Train-Ticket} & 0.032 & 0.016 & 0.008 & \textbf{0.256} \\
\textbf{Sock-Shop-2} & \textbf{0.272} & 0.176 & 0.232 & 0.240 \\
\textbf{Online-Boutique} & \textbf{0.472} & 0.368 & 0.416 & 0.368 \\
\bottomrule
\end{tabular}
\vspace{0.2cm}

\small \textit{註：粗體表示該資料集上的最佳基底函數。PQC\_GPU 在複雜大規模系統（TT）上展現顯著優勢。}
\end{table*}

綜整表 \ref{tab:basis_comparison} 之實證結果，本研究提出以下基底函數選擇策略：

\textbf{策略一：基於系統拓樸結構選擇}
\begin{itemize}
\item \textbf{扇出型拓樸}（如 Online Boutique）：優先選擇 \textbf{Chebyshev} 基底函數，因其穩定逼近特性適合多服務依賴場景。
\item \textbf{線性拓樸}（如 Sock Shop）：優先選擇 \textbf{B-spline} 基底函數，因其分段平滑特性適合線性依賴結構。
\item \textbf{深層複雜拓樸}（如 Train Ticket，40+ 服務）：優先選擇 \textbf{PQC\_GPU} 基底函數，實證顯示其於 Train Ticket 資料集中表現最佳（Precision@1=0.256，Avg@5=0.3648），其指數級表示能力適合複雜多跳傳播場景。
\end{itemize}

\textbf{策略二：基於故障類型選擇}
\begin{itemize}
\item \textbf{資源故障}（CPU、記憶體、磁碟）：優先選擇 \textbf{Chebyshev} 基底函數，實驗顯示其在資源故障上達到最佳性能。
\item \textbf{代碼級故障}：根據系統規模選擇，中等規模選擇 \textbf{B-spline}，大規模選擇 \textbf{Chebyshev} 或 \textbf{PQC\_GPU}。
\end{itemize}

\textbf{策略三：基於資料豐富度選擇}
\begin{itemize}
\item \textbf{純指標資料}（RE1 系列）：優先選擇 \textbf{Chebyshev} 或 \textbf{BARO}，統計方法在資料充分時可能獲得優勢。
\item \textbf{多源資料}（RE2 系列，包含 logs、traces）：優先選擇 \textbf{Fourier} 基底函數，因其能捕捉多源資料中的週期性模式。
\end{itemize}

\section{討論}

\subsection{KAN 有助於 RCA 之定性分析與歸納偏差驗證}

本節從定性與定量兩個層面驗證 KAN 在 RCA 中的優勢。首先，透過可視化分析（見圖 \ref{fig:kan_viz}），我們從\textbf{函數形狀（Function Shape）}層面揭示了 GNN\_KAN 性能優勢的本質來源。其次，我們提出一個量化指標來評估可解釋性。

\begin{figure}[t!]
    \centering
    \includegraphics[width=0.9\columnwidth]{figure2_mlp_vs_kan.jpg}
    \caption{\textbf{KAN 與 MLP 學習到的傳播函數視覺化比較（定性分析）。} 紅色實線（KAN, Proposed）展現了符合微服務故障傳播物理特性的平滑 S 型曲線，具備明顯的「死區（Dead Zone）」與「飽和區（Physical Saturation）」特徵，對應系統在低負載下的韌性與高負載下的崩潰行為；藍色虛線（MLP, Baseline）則呈現無規律的雜訊波動，缺乏可解釋的物理結構。此對比有力證實了 KAN 具備優越的歸納偏差（Inductive Bias），能內化微服務系統的本質特性。}
    \label{fig:kan_viz}
\end{figure}此定性分析提供了超越數值指標的深刻洞察，證實 \textbf{GNN\_KAN 的優勢源於其架構具備符合物理系統特性的歸納偏差，而非僅是參數量的增加}。

\textbf{KAN 學習到的物理語義函數}：在異常傳播路徑的關鍵節點上，KAN 學習到的基底函數呈現出明顯的\textbf{「死區（Dead Zone）」與「飽和區（Saturation Zone）」}結構。具體而言，在圖 \ref{fig:kan_viz} 的紅色實線中，\textbf{死區}對應函數曲線在低輸入值（例如 CPU 使用率 < 60\%）時的平坦區域，表示系統在低負載下的韌性（輸入波動不會觸發異常傳播）；\textbf{飽和區}則對應函數曲線在高輸入值（例如 CPU 使用率 > 90\%）時的漸近線行為，精確捕捉了系統在高負載下的崩潰行為（輸出不再隨輸入線性增長，系統已達到容量上限）。這種\textbf{門檻觸發（Threshold Triggering）}與\textbf{飽和行為（Saturation Behavior）}與微服務系統的實際物理特性高度吻合：當上游服務的延遲或錯誤率超過某閾值時，下游服務開始受到影響（門檻效應）；當異常強度進一步增加時，影響程度趨於飽和，因為系統已達到其處理能力的極限。這證實了 KAN 確實內化了領域知識，而非僅進行黑箱擬合。

\textbf{MLP 的無規律雜訊}：反觀傳統 MLP（使用 ReLU 啟動），其在相同傳播路徑上學習到的函數形狀呈現無規律的雜訊分佈，未能呈現具備物理意義的結構。此證實黑箱 MLP 僅進行資料驅動的擬合，未能捕捉異常傳播的內在規律。實驗結果進一步支持此結論：在複雜場景（Train-Ticket）中，GNN (MLP) 的 P@1 僅為 0.024，幾乎無法定位根因，而 GNN\_KAN 達到 0.256，提升超過 10 倍。這證實了 KAN 的歸納偏差在複雜非線性傳播中的關鍵作用。

\textbf{可視化驗證的方法論意義}：此定性分析有力地回答了「為什麼要用 KAN 而非 MLP」的核心問題。它證實了 GNN\_KAN 的優勢來自\textbf{結構化歸納偏差（Structured Inductive Bias）}，而非單純的模型容量（Model Capacity）。這也是 GNN\_KAN 在有限資料（異常事件罕見）下依然穩健的關鍵原因：KAN 的函數分解限制了函數空間，降低了有效參數維度，避免了過度擬合。

\textbf{基底函數的領域適配性}：透過選擇適合 RCA 的基底函數（B-spline 建模平滑傳播、Chebyshev 穩定逼近、PQC\_GPU 處理超高維交互），我們將領域知識編碼至模型結構中，這是傳統黑箱模型無法實現的設計自由度。

\textbf{可解釋性的量化評估}：為補充定性分析，本研究提出一個量化指標來評估模型解釋與\textbf{結構先驗所誘導之近似傳播路徑}的一致性。對於每個異常事件，我們計算 KAN 學習到的邊權重（反映傳播強度）與已知服務依賴關係所誘導之可達性子圖（reachability subgraph）的交集比例（Intersection over Union, IoU）。具體而言，對於根因節點 $v^*$，我們提取從 $v^*$ 到所有受影響節點的路徑上，KAN 賦予高權重（> 0.7）的邊集合 $E_{\text{KAN}}$，以及在靜態依賴圖上由 $v^*$ 出發、經由所有可達節點所形成的邊集合 $E_{\text{true}}$（作為真實傳播路徑的近似上界）。IoU 定義為：
\begin{equation}
\text{IoU} = \frac{|E_{\text{KAN}} \cap E_{\text{true}}|}{|E_{\text{KAN}} \cup E_{\text{true}}|}
\end{equation}
在 RE1-TT 資料集上的實驗結果顯示，GNN\_KAN 的平均 IoU 達到 0.68，而使用 MLP 的對比方法僅為 0.42。這表明 KAN 學習到的傳播路徑與真實依賴關係具有更高的重合度，從定量角度證實了其可解釋性優勢。然而，我們也承認此指標的局限性：它依賴於完整的服務依賴圖標註，在實務中可能難以獲得。未來工作可探索基於維運專家評估的使用者研究，以更全面地驗證可解釋性。

\subsection{研究限制}

本研究存在以下限制：（1）\textbf{拓樸依賴性}：GNN\_KAN 依賴服務依賴圖作為結構先驗，實務中圖可能不完整或不準確（如追蹤資料遺漏、動態路由），恐影響傳播路徑之捕捉；（2）\textbf{標註偏差}：根因標註由人工事後分析產生，可能包含錯誤或主觀判斷；（3）\textbf{計算開銷}：KAN 的計算開銷略高於 MLP（推論時間增加約 50\%），在極度延遲敏感場景可能需進一步最佳化；（4）\textbf{資料集代表性}：實驗使用的微服務基準雖具代表性，但可能無法涵蓋所有實務場景（如超大規模系統、特定領域應用）；（5）\textbf{絕對性能界線}：在 Train-Ticket 等超大規模複雜場景中，即便 GNN\_KAN 相較基線方法有顯著相對提升（P@1=0.256 vs BARO 0.136 vs GNN 0.024），其絕對 P@1 仍然偏低，反映出在真實多跳級聯與標註不完全的設定下，精確 Top-1 根因定位本身是一個極具挑戰性的開放問題，未來需結合人機協作與 Top-k/路徑級解釋以提升實務可用性。

\section{結論與未來工作}

\subsection{研究總結與展望}

本研究提出 GNN\_KAN，打破根因分析領域長期存在的「表達力與可解釋性兩難」。透過 Kolmogorov–Arnold 表示定理將多變量函數分解為一維可學習函數組合，GNN\_KAN 賦予模型白箱級可解釋性：可視化的「死區/飽和區」函數形狀精確對應系統物理特性。RCA-aware 特徵工程系統化消除後驗洩漏，確保模型學習通用傳播模式而非訓練集偽相關。在多個核心資料集的全面評估中，GNN\_KAN 相較統計方法與傳統 GNN (MLP) 均展現穩定優勢，特別是在超大規模複雜系統中，PQC\_GPU 基底函數顯著優於統計與圖分析基線（Train Ticket: Precision@1=0.256 vs BARO 0.136 vs GNN 0.024 vs PC\_PageRank 0.110），證實其指數級表示能力與結構化歸納偏差在複雜多跳傳播中的優勢。本研究提出系統化基底函數選擇策略，根據拓樸/故障類型/資料豐富度為實務應用提供明確指導。\textbf{GNN\_KAN 標誌根因分析從黑箱邁向可解釋、可信任 AI 的重要一步}，為構建人機協作的下一代智慧維運平台奠定基石。未來可探索：因果推論深度融合、可解釋性量化標準化、AutoML 驅動 KAN 設計、異構圖學習、時序 GNN、少樣本/遷移學習、持續學習、不確定性量化與主動學習等方向。

\begin{thebibliography}{99}
\small

\bibitem{kan2024}
Liu, Z., Wang, Y., Vaidya, S., Ruehle, F., Halverson, J., Soljačić, M., Hou, T. Y., \& Tegmark, M. (2024).
\textit{KAN: Kolmogorov-Arnold Networks}.
arXiv preprint arXiv:2404.19756.

\bibitem{rcaeval2024}
Li, Z., et al. (2024).
\textit{Robust and Explainable Root Cause Analysis for Microservices: A Benchmark}.
In \textit{Proceedings of the ACM Web Conference (WWW)}.

\bibitem{baro2021}
Wu, L., Tordsson, J., Elmroth, E., \& Kao, O. (2021).
\textit{MicroRCA: Root Cause Localization of Performance Issues in Microservices}.
In \textit{Proceedings of the IEEE/IFIP Network Operations and Management Symposium (NOMS)}.

\bibitem{gnn_rca2023}
Wang, X., et al. (2023).
\textit{Graph Neural Network-based Root Cause Analysis for Cloud-Native Systems}.
In \textit{Proceedings of the AAAI Conference on Artificial Intelligence}.

\bibitem{train_ticket}
Zhou, X., et al. (2018).
\textit{Fault Analysis and Debugging of Microservice Systems: Industrial Survey, Benchmark System, and Empirical Study}.
\textit{IEEE Transactions on Software Engineering}, 47(2), 243-260.

\bibitem{sock_shop}
Weaveworks. (2016).
\textit{Sock Shop: A Microservice Demo Application}.
Available: https://microservices-demo.github.io/

\bibitem{online_boutique}
Google Cloud. (2020).
\textit{Online Boutique: A Cloud-Native Microservices Demo Application}.
Available: https://github.com/GoogleCloudPlatform/microservices-demo

\bibitem{pennylane}
Bergholm, V., et al. (2018).
\textit{PennyLane: Automatic differentiation of hybrid quantum-classical computations}.
arXiv preprint arXiv:1811.04968.

\bibitem{kolmogorov1957}
Kolmogorov, A. N. (1957).
\textit{On the representation of continuous functions of several variables by superposition of continuous functions of one variable and addition}.
\textit{Doklady Akademii Nauk SSSR}, 114(5), 953-956.

\bibitem{microservices_survey}
Soldani, J., Tamburri, D. A., \& Van Den Heuvel, W. J. (2018).
\textit{The pains and gains of microservices: A Systematic grey literature review}.
\textit{Journal of Systems and Software}, 146, 215-232.

\bibitem{baro2023}
Li, Z., Zhao, N., Li, M., Ioannidis, S., \& Chowdhury, K. (2023).
\textit{BARO: Robust Root Cause Analysis for Microservices via Multivariate Bayesian Online Change Point Detection}.
In \textit{Proceedings of the ACM Web Conference (WWW)}, 2865-2875.

\bibitem{microrank2021}
Yu, G., et al. (2021).
\textit{MicroRank: End-to-End Latency Issue Localization with Extended Spectrum Analysis in Microservice Environments}.
In \textit{Proceedings of the ACM Web Conference (WWW)}, 3087-3098.

\bibitem{cloudranger2018}
Wang, P., et al. (2018).
\textit{CloudRanger: Root Cause Identification for Cloud Native Systems}.
In \textit{Proceedings of the IEEE/ACM International Symposium on Cluster, Cloud and Grid Computing (CCGRID)}, 492-502.

\bibitem{nezha2023}
Yu, G., et al. (2023).
\textit{Nezha: Interpretable Fine-Grained Root Causes Analysis for Microservices on Multi-modal Observability Data}.
In \textit{Proceedings of the ACM Joint European Software Engineering Conference and Symposium on the Foundations of Software Engineering (ESEC/FSE)}, 428-440.

\bibitem{causalrca2023}
Xin, R., et al. (2023).
\textit{CausalRCA: Causal Inference Based Precise Fine-Grained Root Cause Localization for Microservice Applications}.
\textit{Journal of Systems and Software}, 203, 111724.

\bibitem{xpert2024}
Jiang, Y., et al. (2024).
\textit{Xpert: Empowering Incident Management with Query Recommendations via Large Language Models}.
In \textit{Proceedings of the IEEE/ACM International Conference on Software Engineering (ICSE)}, 1234-1245.

\bibitem{trank2021}
Ye, Z., et al. (2021).
\textit{T-Rank: A Lightweight Spectrum Based Fault Localization Approach for Microservice Systems}.
In \textit{Proceedings of the IEEE/ACM International Symposium on Cluster, Cloud and Grid Computing (CCGRID)}, 61-70.

\bibitem{eadro2023}
Lee, C., et al. (2023).
\textit{EADRO: An End-to-End Troubleshooting Framework for Microservices on Multi-source Data}.
In \textit{Proceedings of the IEEE/ACM International Conference on Software Engineering (ICSE)}, 756-768.

\bibitem{rcacopilot2024}
Chen, Z., et al. (2024).
\textit{RCACopilot: Large Language Model Powered Root Cause Analysis Assistance}.
In \textit{Proceedings of the IEEE/ACM International Conference on Software Engineering (ICSE)}, 2123-2135.

\end{thebibliography}

\end{document}


